\section{The literal activated range}

Fix a scale \(X\), a nonzero shift \(h_0\), and every external choice before
the coefficient vector. Let
\[
 Z_X=\mathbb C^{\mathcal M_X},\qquad
 z=(z_m)_{m\in\mathcal M_X},
\]
be the physical row space. The terminal construction of
\citet{WangTPC64} gives a finite Hilbert space
\[
 \mathcal H_X^{\rm term}
 =\bigoplus_{k\in K_X}^{\perp}
   \ell^2(\mathcal U^+_{X,h_0,k})
\]
and one literal analysis map
\[
 L_X:Z_X\longrightarrow\mathcal H_X^{\rm term}.
\]
Its columns are
\[
 f_m=L_Xe_m,\qquad T_m=\|f_m\|^2.
\]
Different physical rows populate disjoint labelled terminal coordinates at
this stage, so
\begin{equation}\label{eq:terminal-diagonal}
 T_X:=L_X^*L_X=\operatorname{diag}(T_m),\qquad
 \Tform_X(z)=\|L_Xz\|^2=\sum_mT_m|z_m|^2.
\end{equation}
The \emph{activated range} is
\[
 Y_X:=\Ran L_X\subseteq\mathcal H_X^{\rm term}.
\]
It is this range, not the free ambient terminal space, that enters every
represented-reference assertion below.

TPC-64 supplies two compatible parent forms. Write
\[
 R_{X,\mathrm{hyp}}=\operatorname{diag}(P_m),\qquad
 R_{X,\mathrm{rep}}=\operatorname{diag}(\widetilde P_m),
\]
for the literal-hypergraph and canonical-repair forms. Thus
\[
 \mathcal P_{X,\nu}(z)=\langle R_{X,\nu}z,z\rangle,
 \qquad \nu\in\{\mathrm{hyp},\mathrm{rep}\}.
\]
When the finite inflation constant \(\mathfrak I_X\) is available,
\begin{equation}\label{eq:parent-order}
 R_{X,\mathrm{hyp}}\preceq R_{X,\mathrm{rep}}
 \preceq \mathfrak I_XR_{X,\mathrm{hyp}}.
\end{equation}

\begin{definition}[Descending map]
A row-space map \(F_X:Z_X\to H_X\) \emph{descends through terminalization} if
there is a linear map \(A_X:Y_X\to H_X\) such that
\[
 F_X=A_XL_X.
\]
Equivalently, \(\Ker L_X\subseteq\Ker F_X\).
\end{definition}

\begin{lemma}[Exact quotient/range criterion]\label{lem:descent}
A linear map \(F_X:Z_X\to H_X\) descends through \(L_X\) if and only if
\(\Ker L_X\subseteq\Ker F_X\). When it descends, the induced map on \(Y_X\)
is unique.
\end{lemma}

\begin{proof}
Necessity is immediate. Conversely, define \(A_X(L_Xz)=F_Xz\). If
\(L_Xz=L_Xz'\), then \(z-z'\in\Ker L_X\), so \(F_Xz=F_Xz'\). The definition
is well posed, linear, and unique because \(L_XZ_X=Y_X\).
\end{proof}

The distinction is essential. A map acting before terminalization may read
repair-only parent coordinates which terminalize to zero. Such a map need not
descend, and neither the activated-range Gram nor the repair invariance proved
later applies to it.

\subsection{Activation kernels and transported parent metrics}
Let
\[
 S_T=(\Ker T_X)^\perp
     =\operatorname{span}\{e_m:T_m>0\}.
\]
The restriction \(L_X|_{S_T}:S_T\to Y_X\) is an isomorphism. If, for a parent
form \(R_{X,\nu}\),
\begin{equation}\label{eq:kernel-equality}
 \Ker R_{X,\nu}=\Ker T_X,
\end{equation}
then the parent metric transports to \(Y_X\) by
\[
 \widehat R_{X,\nu}
 =(L_X|_{S_T}^{-1})^*R_{X,\nu}|_{S_T}(L_X|_{S_T}^{-1}).
\]
In the normalized row basis \(u_m=f_m/\sqrt{T_m}\), this metric is diagonal:
\begin{equation}\label{eq:transported-parent}
 [\widehat R_{X,\nu}]_{u_m}
 =\operatorname{diag}\left(
 \frac{R_{X,\nu}(m,m)}{T_m}\right)_{T_m>0}.
\end{equation}
This is a transported \emph{parent} metric, not the represented Gram.

\begin{proposition}[Mandatory activation kill]\label{prop:activation-kill}
If there is \(z\in\Ker L_X\) with
\(\langle R_{X,\nu}z,z\rangle>0\), then for every downstream map \(A_X\)
defined on \(Y_X\),
\[
 \inf_{\mathcal P_{X,\nu}(w)>0}
 \frac{\|A_XL_Xw\|^2}{\mathcal P_{X,\nu}(w)}=0.
\]
For the diagonal forms this occurs exactly when some row has
\(R_{X,\nu}(m,m)>0=T_m\).
\end{proposition}

\begin{proof}
Use \(w=z\). The numerator vanishes and the denominator is positive. The
diagonal characterization follows from \eqref{eq:terminal-diagonal}.
\end{proof}

Thus the kernel equality \eqref{eq:kernel-equality} is not cosmetic. Failure
of the reverse inclusion is an activation-layer obstruction that no later
represented operator can repair.
